\documentclass[../../../include/open-logic-section]{subfiles}

\begin{document}

\olfileid{sth}{choice}{hartogs}
\olsection{పోల్చదగినతనం, హార్టోగ్స్ లెమ్మా}

ఇది అనుకూలమైన వైపు. ఇప్పుడు ప్రతికూలమైన వైపు చూద్దాం. ఎంపిక లేకపోతే విషయాలు \emph{సంక్లిష్టం} అవుతాయి. ఎందుకో తెలుసుకోవడానికి \cite{Hartogs1915} ఇచ్చిన ఈ మంచి ఫలితాన్ని చూడండి:

\begin{lem}[\emph{$\ZF$లో}]\ollabel{HartogsLemma}
ఏ సమితి $A$కైనా $\cardnless{\alpha}{A}$ అయ్యే క్రమసంఖ్య $\alpha$ ఉంది.
\end{lem}

\begin{proof}
$B \subseteq A$, $R \subseteq B^2$ అయితే \olref[ord-arithmetic][using-addition]{rankcomputation} ప్రకారం $\tuple{B, R} \subseteq
V_{\setrank{A}+4}$. కాబట్టి విభజనను వాడి దీన్ని తీసుకుందాం:
\[
	C = \Setabs{\tuple{B, R} \in V_{\setrank{A}+5}}{B\subseteq A 
	\text{ మరియు $\tuple{B, R}$ సుక్రమం}}
\]
ప్రతిస్థాపన, \olref[ordinals][ordtype]{thmOrdinalRepresentation}లను వాడి ఈ సమితిని నిర్మిద్దాం:
\[
	\alpha = \Setabs{\ordtype{B, R}}{\tuple{B, R} \in C}.
\]
\olref[ordinals][basic]{corordtransitiveord} ప్రకారం $\alpha$ క్రమసంఖ్య; ఎందుకంటే అది క్రమసంఖ్యల సంక్రమణ సమితి. నిజానికి $\gamma \in \beta \in \alpha$ అయితే, ఏదో $B \subseteq A$కు $\beta = \ordtype{B, R}$. అప్పుడు \olref[ordinals][iso]{wellordinitialsegment} ప్రకారం ఏదో $b
\in B$కు $\gamma = \ordtype{B_b, R_b}$; కాబట్టి $\gamma \in \alpha$.
\sourcecorrection{OLTESTCHOICEHART-001}{మూలంలో సుక్రమిత ఉపసమితి $B$ను $R$కు ఉపసమితిగా చెప్పింది; $B$ అనేది $A$కు ఉపసమితి. ఆ రక లోపాన్ని సరిచేశాం.}

వైరుధ్యం కోసం $f \colon \alpha \to A$ అనే \tetoken{ఇంజెక్షన్}{!!a{injection}} ఉందనుకుందాం. అప్పుడు:
\begin{align*}
	B &= \ran{f}\\
	R &= \Setabs{\tuple{f(\gamma), f(\delta)} \in A \times A}{\gamma \in \delta}.
\end{align*}
స్పష్టంగా $\alpha = \ordtype{B, R}$, $\tuple{B, R} \in C$. అందువల్ల $\alpha
\in \alpha$; ఇది వైరుధ్యం.
\sourcecorrection{OLTESTCHOICEHART-002}{మూలం $R$ నిర్వచనంలో పూర్తి క్రమసంఖ్య $\alpha$నే ప్రమేయ వాదనగా, అంతకంటే చిన్న $\beta$లో సభ్యుడిగా వాడింది. నిర్వచన సమితిలోని రెండు సాధారణ సూచికలను తీసుకుని వాటి క్రమాన్ని చిత్రించిన సంబంధంగా సరిచేశాం.}
\end{proof}

దీని నుంచి లోతైన ఫలితం వస్తుంది:

\begin{thm}[\emph{$\ZF$లో}]
ఈ కింది వాదనలు సమానార్థకాలు:
\begin{enumerate}
	\item\ollabel{equivwo} సుక్రమపరచడం స్వయంసిద్ధం;
	\item\ollabel{equivcompare} ఏ సమితులు $A$, $B$కైనా $\cardle{A}{B}$ లేదా $\cardle{B}{A}$.
\end{enumerate}
\end{thm}

\begin{proof}
\emph{\olref{equivwo} $\Rightarrow$ \olref{equivcompare}.} $A$, $B$లను స్థిరపరచుకుందాం. \olref{equivwo} ప్రకారం $\tuple{A, R}$, $\tuple{B, S}$ అనే సుక్రమాలు ఉన్నాయి. \olref[ordinals][ordtype]{thmOrdinalRepresentation} ప్రకారం $f \colon
\alpha \to A$, $g \colon \beta \to B$లను ఆయా సుక్రమాలతో సమరూపతలుగా తీసుకుందాం. \olref[sth][ordinals][basic]{ordinalsaresubsets} ప్రకారం $\alpha \subseteq \beta$ లేదా $\beta \subseteq \alpha$. $\alpha \subseteq \beta$ అయితే $\comp{f^{-1}}{g} \colon A \to B$ \tetoken{ఇంజెక్షన్}{!!a{injection}}, కనుక $\cardle{A}{B}$; అలాగే $\beta \subseteq \alpha$ అయితే $\cardle{B}{A}$.
\sourcecorrection{OLTESTCHOICEHART-003}{మూలం సమరూపతల లక్ష్యాలుగా క్రమిత నిర్మాణాలనే వ్రాసింది; ప్రమేయాల లక్ష్య సమితులు వరుసగా $A$, $B$ కావాలి. సమరూపత అర్థాన్ని గద్యంలో నిలిపి రకాలను సరిచేశాం.}

\emph{\olref{equivcompare} $\Rightarrow$ \olref{equivwo}.} $A$ను స్థిరపరచుకుందాం; \olref{HartogsLemma} ప్రకారం $\cardnless{\beta}{A}$ అయ్యే ఏదో క్రమసంఖ్య $\beta$ ఉంది. \olref{equivcompare} ప్రకారం $\cardle{A}{\beta}$. కాబట్టి $f \colon A \to
\beta$ అనే \tetoken{ఇంజెక్షన్}{!!{injection}} ఉంది. ఈ ఇంజెక్షన్‌తో $\Setabs{\tuple{a, b} \in A \times A}{f(a)
\in f(b)}$ అనే క్రమాన్ని నిర్వచించి $A$ మూలకాలను సుక్రమపరచవచ్చు.
\end{proof}
\noindent
వెంటనే వచ్చే పరిణామం ఇది: సుక్రమపరచడం విఫలమైతే, కొన్ని సమితుల పరిమాణాలను \emph{అసలు పోల్చలేం}. అందువల్ల అతిపరిమిత కార్డినల్ అంకగణితం సంక్లిష్టమవుతుంది. ఉదాహరణకు $A$, $B$ అనంత సమితులు అయితే, $A$, $B$లలో పెద్దదైన $M$కు $\cardeq{A \disjointsum B}{M}$, $\cardeq{A \times B}{M}$ అని చెప్పే ఆలోచనను వదులుకోవాలి (\olref[card-arithmetic][simp]{cardplustimesmax} చూడండి). సమస్య సులభమే: $A$, $B$ పరిమాణాలను \emph{పోల్చలేకపోతే}, వాటిలో ఏది పెద్దదని అడగడానికే అర్థం లేదు.
\sourcecorrection{OLTESTCHOICEHART-004}{మూల చివరి ఉదాహరణలో సమసంఖ్యకత్వ సంబంధాన్ని మరో సమసంఖ్యకత్వ సంబంధానికి వాదనగా ఇచ్చిన తప్పు సూత్రం ఉంది. రెండు సమితులు పెద్దదానికి సమసంఖ్యకాలన్న ఉద్దేశాన్ని రెండు సక్రమ వాదనలుగా చూపాం.}

\end{document}
